\documentclass[11pt]{article}
\usepackage[T1]{fontenc}
\usepackage[margin=1in]{geometry}
\usepackage{amsmath,amssymb}
\usepackage{graphicx}
\usepackage{booktabs}
\usepackage{longtable}
\usepackage{array}
\usepackage{xcolor}
\usepackage[colorlinks=true,linkcolor=blue!60!black,urlcolor=blue!60!black,
            citecolor=blue!60!black]{hyperref}
\usepackage[htt]{hyphenat}
\usepackage{microtype}
\emergencystretch=3em
\usepackage{listings}
\lstset{basicstyle=\ttfamily\footnotesize, backgroundcolor=\color{gray!8},
        frame=single, rulecolor=\color{gray!40}, breaklines=true,
        columns=fullflexible, keepspaces=true, showstringspaces=false}

\newcommand{\code}[1]{\texttt{#1}}
\newcommand{\file}[1]{\texttt{\small #1}}
\newcommand{\fn}[1]{\texttt{\small #1}}
\newcommand{\Om}{\Omega_0}
\newcommand{\qom}{q\Om}

\newcommand{\invariant}[2]{%
  \par\vspace{0.6em}\noindent
  \fcolorbox{gray!40}{gray!8}{\parbox{0.97\linewidth}{\textbf{#1}\quad #2}}%
  \par\vspace{0.6em}}

\title{Dust Particles and Gas$+$Dust Self-Gravity\\
       in the AthenaK Shearing Box\\[0.4em]
       \large \textsc{Implementation Document}\\
       \large Phase 4a: the dust-module merge; Phase 4b (i): the shear-periodic
       deposit fold; (ii): the ideal-gas energy update; (iii): the particle
       timestep under shear}
\author{\code{athenak-multigrid} tree, branch \code{dust-multigrid}\\
        \small commits \code{449b8f09} (merge), \code{a7c0def3}, \code{774c6519},
        \code{78320aee} (second merge), \code{6c0cafef}, \code{489c656c}, \code{375643e6}}
\date{September 3--4, 2026}

\begin{document}
\maketitle

\begin{abstract}
\noindent
This document describes \emph{how} Lagrangian dust particles are being brought into
the \code{athenak-multigrid} tree --- the shearing-box code with multigrid
self-gravity on refined meshes documented in
\file{multigrid\_shear\_implementation.pdf} --- on the way to gas$+$dust
self-gravity in a stratified, cooling, gravito-turbulent box (the configuration of
Baehr, Zhu \& Yang 2022, with the radial pressure gradient they omitted). It is the
companion to \file{dust\_selfgravity\_validation.pdf}, which records what the
resulting code does on test problems. Both documents are meant to grow phase by phase.

\textbf{Phase 4a}, recorded here, merges the sibling fork \code{athenak-dust}
(particle-mesh dust with IMEX drag, validated in its own repository) into the
multigrid tree. The merge itself is mechanical --- eight keep-both conflicts and one
renumbered output slot --- but it exposed a real defect: the dust module's additive
ghost-deposit exchange was written against the per-neighbor MPI messaging of
July's upstream, while this tree carries the rank-packed messaging of upstream
\#775. Any multi-rank dust run deadlocked at cycle 0. The exchange now uses the
aggregated-message path with the same idiom as the cell-centered exchange, and
ranks 1, 2 and 4 agree to every printed digit. A second follow-up aligns the dust
task graph with the tree's Prolongate-before-boundary-conditions order.

A \textbf{second merge} the same day caught the tree up with the dust fork's own
latest work (seven commits, to \code{1f3e1ed9}): the fork had independently made
the identical rank-packed rewrite of the deposit exchange, so that conflict resolves
to their version; and it now offers two \emph{predictor--corrector} coupling modes
(\code{coupling = pc2} and the adaptive \code{hybrid}) that run under plain
\code{rk2}, second order in the resolved drag regime with a stiff backward-Euler
fallback. Dust can therefore run under the same integrator as every gravity
validation before it; the \code{imex2+} path remains for stiff cases.

\textbf{Phase 4b (i)} closes the first of the gaps the survey named: ghost
deposits of the MeshBlocks on the shear-periodic $x_1$ faces are now folded across the
faces with a conservative azimuthal remap --- gathered into global $y$--$z$ planes,
shifted by the shear offset with the Phase-1 remap-flux kernels, and \emph{added} into
the active edge strips of the blocks across the face --- while the unsheared
plain-periodic $x_1$ contributions of those blocks are suppressed. The fold is the
adjoint of the copy-exchange ghost fill, so its shift signs are the reverse of the
multigrid fill's. A unit-test problem generator certifies it exactly at integer shifts
and at second/third order otherwise, and a structured 3D run shows the old unsheared
exchange was a percent-level error. \textbf{Phase 4b (ii)} lifts the isothermal
restriction on back-reaction: every drag kick on the gas momentum now goes through one
helper that, for an ideal gas, books the matching kinetic-energy change so the internal
energy is untouched, and a flagged option deposits the frictional dissipation of each
particle kick as heat. The residual energy errors are the convexity gap of the
low-storage RK combination, first order in the timestep, and are quantified.
\textbf{Phase 4b (iii)} removes the last of the three Phase-4b gaps: the particle
transport timestep no longer resolves the background shear $-\qom x$ cell by cell (a
$1/L_x$ penalty the orbitally advected gas does not pay) but limits the peculiar
velocity per cell and the \emph{total} transport velocity per MeshBlock --- the latter
being what the single-hop particle migration and the one-layer deposit halo actually
require --- with the bound made exact for the low-storage tableaus through their largest
stage fraction. In an $L_x = 8$ box the step grows $5.3\times$ to the gas limit.

The document also records, for the phases that follow, the integration map derived
from the code survey: where the gravity solve is called relative to the dust task
graph, the three places the solver reads density, the ghost-width contract of the
potential, and the design chosen for the particle force (grid gradient, exchanged
and shear-remapped like the module's provisional gas velocity, gathered with the
deposit kernel), together with the known gaps that Phases 4b--4d must close.
\end{abstract}

\tableofcontents
\newpage

%=========================================================================================
\section{Scope, sources, and how to read this}
\label{sec:scope}

\subsection{What this document is}

This is an implementation reference for the dust track. For each phase it answers:
what the code relies on, where that is established and enforced, which data
structure carries the new state, which function does the work, and from which call
site it is reached. Results live in the companion documents:

\begin{center}
\small
\begin{tabular}{@{}>{\raggedright\arraybackslash}p{7.4cm}>{\raggedright\arraybackslash}p{7.6cm}@{}}
\toprule
\file{validation/dust\_selfgravity\_validation.pdf} & results, tests, figures (this track) \\
\file{validation/multigrid\_shear\_implementation.pdf} & the gravity/refinement machinery being extended \\
\file{validation/multigrid\_selfgravity\_validation.pdf} & its validation record (Phases 0--3) \\
\file{../athenak\_dust\_tests/} & the dust module's own validation notebooks and inputs \\
\bottomrule
\end{tabular}
\end{center}

\subsection{The program}
\label{sec:program}

The target science is the interaction of gravitational instability of the gas with
dust concentration and collapse in a shearing box: Baehr, Zhu \& Yang (2022, ApJ 933,
100) in 3D stratified, cooling, gravito-turbulent boxes with superparticles at
$\mathrm{St} = 0.1$--$10$, plus the radial pressure gradient (and hence the streaming
instability) that their setup left out, and eventually adaptive zoom on the dust
clumps. The phases:

\begin{center}
\small
\begin{tabular}{@{}llp{8.6cm}@{}}
\toprule
Phase & Status & Content \\
\midrule
4a & \textbf{done} & merge \code{athenak-dust} into the multigrid tree; establish that both
     the dust module and the gravity machinery are unchanged by the merge \\
4b & (i), (ii) \textbf{done} & (i) shear-periodic \emph{additive} deposit remap
     (\S\ref{sec:phase4b1}); (ii) ideal-gas energy update under drag back-reaction
     (\S\ref{sec:phase4b2}); (iii) particle timestep on the shear-subtracted velocity;
     adopt the PC2 coupling (\S\ref{sec:merge2}) as the production path \\
4c & planned & dust self-gravity: mass deposit into the Poisson source, three density
     reads, force gather to particles; dust-only and gas$+$dust tests \\
4d & planned & the science configuration: SC14 box $+$ pressure gradient $+$ dust;
     particle restart; clump diagnostics \\
5  & deferred & particles on refined meshes (deposit exchange across levels, particle
     redistribution on remesh) \\
\bottomrule
\end{tabular}
\end{center}

\subsection{Commit provenance}

\begin{table}[htb]
\centering
\small
\begin{tabular}{@{}llp{9.4cm}@{}}
\toprule
Commit & Phase & Content \\
\midrule
\code{449b8f09} & 4a & Merge of \code{dust/dust} (\code{d6de1296}) into \code{multigrid}
                       (\code{b4d4c547}); eight conflicts resolved keep-both; output slot
                       \code{grav\_phi} $153 \to 154$ \\
\code{a7c0def3} & 4a & Dust task graph: \fn{Prolongate} before \fn{ApplyPhysicalBCs} \\
\code{774c6519} & 4a & \file{bvals\_dep.cpp}: additive deposit exchange on the rank-packed
                       MPI vars path (multi-rank deadlock fix) \\
\code{78320aee} & 4a & Second merge, \code{dust/dust} at \code{1f3e1ed9}: the fork's own
                       rank-packed deposit exchange (taken), PC2/hybrid coupling, PM
                       optimizations; three conflicts (\S\ref{sec:merge2}) \\
\code{6c0cafef} & 4b (i) & \fn{MeshBoundaryValuesDep::FoldShearDeposit}, the $x_1$-buffer skip,
                  the dust hooks, the \code{dust\_deposit\_shear} test generator, the
                  3D shear-periodic NSH input, the NSH mass modulation \\
\code{489c656c} & 4b (ii) & \fn{GasKick}, the four kick sites, the heat channel of the
                  back-reaction deposit, \code{<dust>/drag\_heating}, the isothermal
                  guard removed, ideal-gas initialization of \code{dust\_nsh} \\
\code{375643e6} & 4b (iii) & \fn{DustGasDrag::ComputeNewTimeStep}: cell limit on the peculiar
                  velocity, MeshBlock guard on the total velocity;
                  \code{<dust>/dt\_transport}, \code{dt\_block\_safety};
                  \code{DUST\_DT\_SUMMARY} \\
\bottomrule
\end{tabular}
\caption{Implementation commits covered by this document.}
\label{tab:commits}
\end{table}

The merged fork is \code{github.com/athenak-dust/athenak}, branch \code{dust}. At
the first merge it stood at \code{d6de1296}, thirteen commits on top of upstream
AthenaK as of July 8, 2026 (commit \code{5f199310}, the merge base with this tree).
The one that matters is \code{e785735e} ``Add IMEX dust-particle coupling to
hydrodynamics'' (36 files, $+3864$ lines); the rest are the matrix-free coupled drag
solvers, GPU migration fixes, a snapshot-seeded clump benchmark, and diagnostics.
The local clone turned out to be two weeks stale; the second merge (\S\ref{sec:merge2})
brought in the seven commits through \code{1f3e1ed9} (September 2). The remote is
registered locally as \code{dust} (\code{git remote add dust ../athenak-dust}), so
later fixes in that fork come in with a plain \code{git merge dust/dust} --- which is
exactly what the second merge was.

\subsection{Files touched}

\begin{center}
\small
\begin{tabular}{@{}>{\raggedright\arraybackslash}p{6.6cm}>{\raggedright\arraybackslash}p{9.2cm}@{}}
\toprule
File & Role \\
\midrule
\file{src/dust/dust.\{cpp,hpp\}} & \code{DustGasDrag} module: parameters, arrays,
  boundary objects, guards, the shared PM kernel \fn{dust::PMWeights} \\
\file{src/dust/dust\_tasks.cpp} & the combined hydro$+$particle task graph
  (\fn{AssembleDustGasDragTasks}) and the boundary tasks \\
\file{src/dust/dust\_pm.cpp} & deposit, per-cell implicit solve, gather/kick, back-reaction \\
\file{src/dust/dust\_push.cpp} & explicit particle push (Coriolis, tide, drift, drag history) \\
\file{src/dust/dust\_newdt.cpp} & particle transport timestep, stopping-time bookkeeping,
  the $\gamma$-switch \\
\file{src/dust/dust\_solver.cpp} & optional matrix-free coupled drag solvers \\
\file{src/bvals/bvals\_dep.cpp} & \code{MeshBoundaryValuesDep}: additive ghost-deposit
  exchange (rewired in 4a) \\
\file{src/bvals/bvals\_part.cpp}, \file{bvals.\{cpp,hpp\}} & shear-periodic particle
  routing; deposit-exchange class declaration \\
\file{src/driver/driver.\{cpp,hpp\}} & \fn{SetImEx2PlusCoefficients} (the $\gamma$-switch
  entry point) \\
\file{src/mesh/meshblock\_pack.\{cpp,hpp\}} & \code{pdust} registration; hydro/particle
  task lists ceded to the dust module when a \code{<dust>} block exists \\
\file{src/mesh/mesh.\{cpp,hpp\}}, \file{build\_tree.cpp} & \code{<time>/dtmax} hard cap \\
\file{src/athena.hpp} & widened \code{ParticlesIndex} (17 real, 3 integer fields) \\
\file{src/outputs/*} & \code{dust\_d} derived output; particle velocities in \code{pvtk};
  slot renumbering \\
\file{src/pgen/tests/dust\_*.cpp}, \file{streaming\_linear.cpp} & the four built-in dust
  problem generators \\
\file{inputs/dust/*.athinput} & their reference inputs \\
\bottomrule
\end{tabular}
\end{center}

%=========================================================================================
\section{The target problem and the design it implies}
\label{sec:design}

\subsection{What is being built}

The reference implementation for particle self-gravity in a shearing box is the
Athena-C problem generator \file{parsg\_strat3d.c} in \file{athena\_parsg/src/prob/}
(Rixin Li's fork). Its method, with \code{PLUS\_GAS\_SELF\_GRAVITY}: deposit the dust mass onto the grid with
the TSC kernel, \emph{add the gas density into the same buffer}, solve \emph{one}
Poisson equation, difference $\Phi$ on the grid, and interpolate the resulting
acceleration to the particles with the \emph{same} TSC kernel used for the deposit.
The gas feels the total potential as a body force. Baehr et al.\ do the same in
\textsc{Pencil} (their eq.~15, $\rho = \rho_g + \rho_d$ in a periodic FFT). Two
details of the Athena-C version are not to be copied: its deposit omits the
cell-volume division and hides it in the particle mass, and its drag feedback must
subtract the gravitational kick because gravity is folded into the same implicit
particle update as drag.

\subsection{The integration map}
\label{sec:map}

The survey of this tree fixed where each piece attaches. These are facts about the
code as merged; the items marked \emph{planned} are the Phase 4b/4c designs.

\paragraph{Where gravity is solved.} Gravity has no task-list entry. The driver
calls \fn{Gravity::Solve(stage)} once per explicit stage, \emph{between} the
\code{before\_stagen} and \code{stagen} task lists
(\file{driver/driver.cpp}, lines 563--569). Anything the solver must see --- the dust mass
deposit and its ghost exchange --- has to complete inside \code{before\_stagen}
(\emph{planned}: a deposit task there, on the stage-start particle positions).

\paragraph{The three density reads.} The solver reads gas density in three places,
and all three must learn about dust (\emph{planned}, Phase 4c):
\begin{itemize}\itemsep2pt
  \item \file{gravity/mg\_gravity.cpp:353}, \fn{LoadSource(u0, IDN, \ldots)} into the
        multigrid source: loads active zones \emph{plus one ghost layer}, so the dust
        deposit must already be additively ghost-exchanged one layer deep;
  \item \file{multigrid/multigrid\_slab.cpp:263}, the slab-plane gather that builds the
        open-$z$ Dirichlet planes (active zones, reduced across ranks) --- omitting dust
        here would make the planes inconsistent with the interior source;
  \item \file{gravity/fft\_gravity.cpp:218}, the FFT oracle's density gather.
\end{itemize}

\paragraph{The potential's ghost contract.} \code{phi} is allocated with the hydro
ghost width but only \emph{one} ghost layer is ever filled
(\code{mg\_nghost = 1}; larger values are fatal with shear). The gas source term
(\fn{SourceTerms::SelfGravity}, \file{srcterms/srcterms.cpp:298}) needs exactly
that: a centered $-\rho\nabla\Phi$ for momentum and the flux-based energy term of
Mullen, Hanawa \& Gammie (2020). A TSC gather at a particle in the last active cell
reaches one ghost cell of the \emph{acceleration}, i.e.\ two of the potential.
\emph{Planned}: compute $\mathbf g = -\nabla\Phi$ on active cells, then copy-exchange
and shear-remap the three-component field exactly as the module already does for its
provisional gas velocity $u^\ast$ (\code{MeshBoundaryValuesCC} plus a
\code{ShearingBoxCC}), and gather with \fn{dust::PMWeights}. Using the deposit kernel
for the force gather --- not the derivative of the kernel --- is the
Hockney--Eastwood choice that suppresses self-force, and it is what the Athena-C
reference does.

\paragraph{Where the particle force goes.} \fn{DustGasDrag::ExplicitPush}
(\file{dust/dust\_push.cpp:65}) already assembles the Coriolis and vertical-tide
accelerations per particle; the gathered $\mathbf g$ joins them there. The
back-reaction deposits only the drag kick $\Delta v_j = c_j(\tilde u^\ast_j - v_j)$,
so gravity is never double-counted --- there is no analogue of the Athena-C
\fn{feedback\_corrector} subtraction to write.

\paragraph{The gas force.} Nothing to add: the existing source term applies whatever
potential the solver produced.

\paragraph{Stage structure.} Under \code{coupling = imex}, \code{imex2+} has three
explicit stages of which the third is dormant for explicit physics ($\beta_3 = 0$);
the gravity solve should be skipped there or a third of the solver cost is wasted
(\emph{planned}). Under the PC2/hybrid couplings of \S\ref{sec:merge2} the integrator
is \code{rk2} and the gravity solve runs exactly as it does for gas alone.

%=========================================================================================
\section{The dust module as inherited}
\label{sec:module}

This section is the working summary of \code{athenak-dust}, with the facts the later
phases build on. The design is Yang \& Johansen (2016) particle-mesh coupling (one
shared kernel for scatter and gather, exact momentum-conserving back-reaction) with
the Benítez-Llambay, Krapp \& Pessah (2019) closed-form per-cell implicit drag inside
the Krapp et al.\ (2024) IMEX(4,3,2) integrator (\code{imex2+}), generalized to
per-particle stopping times.

\subsection{Data layout}

Particles are two flat device arrays indexed \code{(field, particle)}. The field
enumeration (\file{athena.hpp:76}) is
\begin{lstlisting}
enum ParticlesIndex {PGID=0, PTAG=1, PSP=2,
                     IPX=0, IPVX=1, IPY=2, IPVY=3, IPZ=4, IPVZ=5,
                     IPX1=6, IPVX1=7, IPY1=8, IPVY1=9, IPZ1=10, IPVZ1=11,
                     IPRX=12, IPRY=13, IPRZ=14, IPTS=15, IPM=16};
\end{lstlisting}
so a dust particle carries position, velocity (all three components even in 2D
$r$--$z$), the low-storage stage-1 registers, the recorded drag rate
$R_j$, its stopping time and its \textbf{mass}. Integer data: owning MeshBlock,
global tag, species. \code{nrdata = 17}, \code{nidata = 3}.

\subsection{The particle-mesh kernel}

\fn{dust::PMWeights} (\file{dust/dust.hpp:78}) returns the three weights of a
NGP / CIC / TSC stencil in one direction; every deposit, gather, matrix-free
operator and kick uses it, which is what makes scatter and gather weights identical
by construction. \code{<dust>/deposit} selects the kernel (default \code{tsc}); the
stencil is always three wide (NGP zeroes the wings). Particle-mesh assignment needs
\code{nghost >= 2}.

\subsection{The stage algorithm}
\label{sec:stage}

Per explicit stage, on the conserved gas state (primitives are stale mid-stage),
with $a\,\Delta t = a_{\rm impl}\Delta t$ the implicit weight of the tableau:
\begin{enumerate}\itemsep2pt
  \item \fn{DepositDrag}: scatter $Q = \sum_j \mu_j W_{jk}$ and
        $\mathbf P = \sum_j \mu_j W_{jk}\mathbf v_j$ with $c_j = a\Delta t/(t_{s,j} + a\Delta t)$,
        $\mu_j = m_j c_j / V$ (a \emph{drag-weighted} density, not a mass density);
        additive halo exchange.
  \item \fn{GasImplicitSolve}: per cell, $\mathbf u^\ast = (\rho\mathbf u + \mathbf P)/(\rho + Q)$,
        the closed-form momentum-conserving backward-Euler pair. With
        \code{back\_reaction = false} the deposits are excluded so test particles are
        kicked toward the unmodified gas.
  \item Copy exchange of $\mathbf u^\ast$ ghosts (plus shear remap in a 3D shearing box).
  \item \fn{GatherKickPMBR}: gather $\tilde{\mathbf u}^\ast_j$, kick
        $\Delta\mathbf v_j = c_j(\tilde{\mathbf u}^\ast_j - \mathbf v_j)$, record
        $R_j = \Delta\mathbf v_j/(a\Delta t)$, and deposit
        $-m_j W_{jk}\Delta\mathbf v_j/V$ with the \emph{same} weights and the same
        $\Delta\mathbf v$ --- exact momentum conservation to round-off.
  \item Additive halo exchange of that deposit; \fn{ApplyPMBR} adds it to the gas
        momentum and rescales it in place into the gas drag rate $R_g$.
  \item Next stage, the history terms $\tilde a_{22}\Delta t\,R$ enter on both sides
        (\fn{AddDragHistoryGas}; the particle side inside \fn{ExplicitPush}).
\end{enumerate}
There is \emph{no} drag timestep constraint for any $t_s$ or dust-to-gas ratio; only
particle transport limits $\Delta t$ (\S\ref{sec:dt}).

\subsection{The task graph}
\label{sec:graph}

When a \code{<dust>} block exists, \fn{MeshBlockPack::AddPhysics} skips both
\fn{Hydro::AssembleHydroTasks} and \fn{Particles::AssembleTasks} and hands the whole
graph to \fn{DustGasDrag::AssembleDustGasDragTasks} (\file{mesh/meshblock\_pack.cpp:265--271},
\file{dust/dust\_tasks.cpp:37--108}), the same pattern as \code{IonNeutral}. The
consequence for later phases:

\invariant{Ownership.}{The dust module owns the entire per-stage task graph. Any
task the gravity work needs (deposit, exchange, force gather) is inserted into
\fn{AssembleDustGasDragTasks}, and the hydro chain inside it must track changes made
to \fn{Hydro::AssembleHydroTasks} in this tree (the 4a follow-up
\code{a7c0def3} is the first instance).}

The \code{stagen} sequence: \fn{FirstTwoImpRK} (replaces \fn{CopyCons}) $\to$ gas
fluxes / update / source terms $\to$ \fn{AddDragHistoryGas}; in parallel
\fn{ExplicitPush} $\to$ particle migration (\fn{NewGID}, counts, sends, receives ---
\emph{inside every stage}, not once per cycle) $\to$ the deposit chain of
\S\ref{sec:stage} $\to$ the gas tail (orbital advection, restrict, exchange, shear
exchange, \fn{Prolongate}, \fn{ApplyPhysicalBCs}, \fn{ConToPrim}, timesteps). Every
dust task returns early on the dormant third stage.

\subsection{Shearing-box conventions}

Particle velocities are \emph{shear-subtracted}: the explicit push applies the
FARGO-form Coriolis terms $f_x = 2\Om v_y$, $f_y = -(2-q)\Om v_x$ and, with
\code{<shearing\_box>/stratified}, $f_z = -\Om^2 z$
(\file{dust/dust\_push.cpp:98--100}); positions drift with the \emph{full}
azimuthal velocity, $\dot y = v_y - \qom x$ (\file{:113--114}). This matches the gas
with \code{orbital\_advection = true} (the tree's default and the SC14 setting).
Crossing a radial face is therefore a \emph{position-only} transform: wrap $x$, shift
$y$ by $\mp y_{\rm sh}$, $y_{\rm sh} = \qom L_x t \bmod L_y$
(\file{bvals/bvals\_part.cpp:76,134}), with the RK position registers shifted by the
same amounts so the low-storage combination stays in one periodic image. The
destination MeshBlock is found from a precomputed $(y,z)\to$ (gid, rank) map of the
face blocks, not from the $x_1$ neighbor.

The radial pressure gradient is a gas-only constant acceleration through the stock
\code{<hydro\_srcterms>/const\_accel} machinery ($2\Om\eta v_K$ along $x$); particles
never feel it.

\subsection{The particle timestep}
\label{sec:dt}

\fn{DustGasDrag::NewTimeStep} (\file{dust/dust\_newdt.cpp:95}) takes
$\min_p \min_i \Delta x_i/|v_{{\rm transport},i}|$ times \code{<dust>/dt\_cfl}
(default $0.5$), where in a 3D shearing box $v_{{\rm transport},y} = v_y - \qom x$
(\file{:120}). In a wide box this is the same penalty the gas pays without orbital
advection ($\sim 5\times$ at SC14 size). Phase 4b (iii) relaxes it
(\S\ref{sec:phase4b3}).

\subsection{Guards and parameters}

Fatal conditions in \fn{DustGasDrag::DustGasDrag} (\file{dust/dust.cpp}): MHD
present (\code{:45}); 1D mesh (\code{:50}); an integrator other than \code{imex2+}
(\code{:63}); \code{multilevel} (\code{:69}); boundaries neither strictly periodic
nor 3D with shear-periodic $x_1$ (\code{:78}); \code{nghost < 2} (\code{:93});
back-reaction with a non-isothermal EOS (\code{:170}); a coupled solver other than
\code{local} with shear (\code{:227}). A \emph{warning} at \code{:84} states that
ghost deposits are exchanged with the plain-periodic pattern under shear, exact only
for azimuthally uniform states. The same \code{multilevel} guard is repeated in the
\code{MeshBoundaryValuesDep} constructor and in the shear-periodic particle
constructor. The parameter interface is reproduced in \S\ref{sec:params}.

%=========================================================================================
\section{Phase 4a: the merge}
\label{sec:phase4a}

\subsection{Route}

The two forks share upstream history through \code{5f199310} (July 8, 2026); the
dust fork adds 13 commits, this tree about 140 including an upstream merge on
August 10 that brought in \#775 (rank-packed boundary messaging) and \#612
(\code{imex2+}). A merge was preferred over cherry-picks: it keeps the dust fork's
history and authorship, resolves conflicts once, and makes later fixes from that
fork a one-line pull. The whole sequence was first rehearsed in a scratch clone
(conflicts, follow-ups, both builds, the full acceptance battery) and then repeated
in the working tree with the rehearsed resolutions copied in.

\subsection{Conflicts}

Eighteen files were touched by both forks; eight conflicted, all of the same kind
--- both sides appended adjacent code at the same anchor --- and all were resolved by
keeping both sides:

\begin{center}
\small
\begin{tabular}{@{}>{\raggedright\arraybackslash}p{4.6cm}>{\raggedright\arraybackslash}p{5.4cm}>{\raggedright\arraybackslash}p{5.8cm}@{}}
\toprule
File & this tree had added & the dust fork had added \\
\midrule
\file{driver/driver.cpp} & the super-time-stepping controller functions &
  \fn{SetImEx2PlusCoefficients} (the constructor's inline \code{imex2+} tableau moved
  into a function the $\gamma$-switch can call) \\
\file{mesh/mesh.cpp}, \file{mesh.hpp} & STS timestep cap and members & \code{dtmax}
  hard cap (default $\infty$, inert) \\
\file{mesh/meshblock\_pack.cpp}, \file{.hpp} & physics (9) gravity, \code{pgrav} &
  physics (10) dust, \code{pdust} \\
\file{outputs/outputs.hpp}, \file{basetype\_output.cpp} & \code{grav\_phi} at slot 153
  and its guard & \code{dust\_d} at slot 153 and its guard \\
\file{parameter\_input.cpp} & \code{gravity} in the block-name list & \code{dust} \\
\bottomrule
\end{tabular}
\end{center}

\paragraph{The output slot.} Output variable names live in one static array,
\code{var\_choice}, and the index of the matching name (\code{ivar}) is used only by
the sanity guards that test numeric ranges (``151--153 are particle outputs, so a
\code{<particles>} block must exist''). Each fork had appended one name at position
153 and written a guard for it. Resolution: \code{dust\_d} keeps 153,
\code{grav\_phi} becomes 154, \code{NOUTPUT\_CHOICES} $154 \to 155$, the gravity
guard tests 154. Everything downstream selects by name string, so the number appears
nowhere else. The 31 auto-merged files were verified identical to the rehearsal.

\subsection{Follow-up 1: task order}

Upstream reordered the hydro chain to \fn{Prolongate} $\to$ \fn{ApplyPhysicalBCs}
$\to$ \fn{ConToPrim} (so physical boundaries act on prolongated ghosts); the dust
graph still had the July order. Commit \code{a7c0def3} aligns
\file{dust/dust\_tasks.cpp:96--98}. Inert on uniform meshes (dust is fatal on refined
ones), but the invariant of \S\ref{sec:graph} says keep them in lockstep.

\subsection{Follow-up 2: the deposit exchange on the rank-packed path}
\label{sec:rankpacked}

\paragraph{Symptom.} The first multi-rank test of the merged tree (the damping
problem with drifting particles, \code{v0 = 0.5}, 4 ranks) hung at cycle 0 with all
four processes spinning; the 1-rank run of the same binary reproduced the reference
number exactly, and multigrid tests at 4 ranks were bit-identical to serial.

\paragraph{Mechanism.} \code{MeshBoundaryValuesDep} (\file{bvals/bvals\_dep.cpp})
was written as a copy of July's \code{MeshBoundaryValuesCC}: per-neighbor
\code{MPI\_Isend} on \code{sendbuf[n].vars\_req[m]} with
\fn{CreateBvals\_MPI\_Tag(lid, dn)} tags, and per-neighbor \code{MPI\_Test} on
\code{recvbuf[n].vars\_req[m]}, relying on the base class \fn{InitRecv} to post the
matching per-neighbor receives. In this tree \fn{MeshBoundaryValues::InitRecv}
(\file{bvals/bvals\_tasks.cpp:32}) is \emph{rank-packed} since upstream \#775: it
builds, once per mesh sequence, a metadata table of all (block, neighbor) payloads
per peer rank (\fn{BuildRankPackedVarMetadata}, \file{bvals/bvals.cpp:134}, which
also does a one-shot header exchange so receivers know the pack order), and posts
\emph{one} \code{MPI\_Irecv} per peer rank with tag 1 into
\code{rank\_recvbuf\_vars\_}. \fn{ClearRecv}/\fn{ClearSend} wait on those aggregated
requests. So the deposit's per-neighbor sends had no receiver, its tests polled
requests that were never posted, and the stage never completed.

\paragraph{Fix.} Commit \code{774c6519} re-expresses the two transport steps in the
same idiom as the cell-centered exchange (\fn{PackAndSendCC} and \fn{RecvAndUnpackCC}
of \code{MeshBoundaryValuesCC}):
\begin{itemize}\itemsep2pt
  \item \fn{PackAndSendDeposit} (\file{:122}): on-rank neighbors are still written
        straight into the destination's \code{recvbuf[dn].vars}; off-rank payloads go
        into \code{rank\_sendbuf\_vars\_} at \code{send\_agg\_offset\_(m*nnghbr+n)}; after
        one fence, one \code{MPI\_Isend} per message in \code{send\_var\_msgs\_} (tag 1,
        \code{send\_var\_reqs\_});
  \item \fn{RecvAndSumDeposit} (\file{:228}): \code{MPI\_Test} over
        \code{recv\_var\_reqs\_}; in the (deliberately scalar) loop over neighbors the
        addend is \code{aggrbuf(base + bi)} when \code{recv\_agg\_offset\_(m*nnghbr+n) >= 0}
        and \code{rbuf[n].vars(m, bi)} otherwise.
\end{itemize}
The metadata builder sizes each entry from the buffer index ranges the derived class
defines (\fn{GetVarDataSize} $\to$ \code{isame\_ndat}), so the reversed data flow of
the additive exchange (pack ghosts, sum into active strips) needs no change to the
base class. The serial path is unchanged (bit-identical damping ladder); ranks 1, 2
and 4 now agree to every printed digit.

\invariant{Messaging.}{Every new boundary-exchange class in this tree must use the
rank-packed vars path (\code{send\_var\_msgs\_} / \code{recv\_var\_reqs\_} /
\code{*\_agg\_offset\_}); the per-neighbor \code{vars\_req} arrays are dead for vars
traffic and only the flux-correction path still posts per-neighbor requests.}

\subsection{The second merge: catching up with the dust fork}
\label{sec:merge2}

\paragraph{What came in.} A question about a colleague's commit revealed that the
local dust clone was two weeks behind. \code{origin/dust} at \code{1f3e1ed9} adds,
oldest first: a CPU particle-mesh optimization (plain addition instead of atomics
under \code{Kokkos::Serial}, weights evaluated in cell coordinates); the NSH generator
without dust; \emph{their own} import of upstream \#775 (\code{0d7688b6}),
extended to the additive deposit exchange --- the same rewrite as \code{774c6519},
validated by them at 2, 4 and 8 ranks and on CUDA; adaptive PC2 / split-BE coupling
(\code{4b44ba09}); and direct PC2 (\code{7fc41cd8}).

\paragraph{Conflicts.} Three files. \file{bvals/bvals\_dep.cpp}: both sides rewrote
the same functions the same way; the fork's version was taken (it additionally
rebuilds the rank-packed metadata inside \fn{PackAndSendDeposit} when the variable
count or mesh sequence changed, instead of relying on \fn{InitRecv} having run).
\file{bvals/bvals\_fc.cpp}: ours kept --- the early return on inactive face-centered
faces is upstream \#739's $\nabla\!\cdot\!B$ fix, which the fork's \#775 import predates.
\file{mesh/mesh\_refinement.cpp}: ours kept --- the fork calls
\fn{InitBoundaryValuesAndPrimitives} at the ancestor's position, while the Phase-0c
code here calls it later, after the shear-state rebuild; keeping both would call it
twice. The base boundary-value files auto-merged with \emph{no net change}: the
fork's \#775 import is textually the upstream copy this tree already had.

\paragraph{The PC2 and hybrid couplings.} \code{<dust>/coupling} now selects
\code{imex} (the default, \S\ref{sec:stage}, requires \code{imex2+}), \code{hybrid}, or
\code{pc2} (both require \code{rk2}; \code{gamma\_switch} is ignored with a warning).
From the fork's own description: the \emph{PC2} branch advances the particles once
per cycle without a coupled gas solve --- a fresh old-state drag-feedback predictor
is deposited during RK stage 1, the midpoint gas velocity is built after the completed
stage-1 boundary tail, the particles take an implicit-midpoint drag update at the
predicted midpoint position, the actual drag impulse is deposited at that midpoint,
and the final gas momentum is corrected by removing the tableau-weighted predictor.
Particles migrate once, after the update. The \emph{split-BE} branch first completes a
dust-free hydro RK2 step and then performs one full-$\Delta t$ coupled backward-Euler
drag solve (first order, L-stable). \code{hybrid} switches between them globally with
hysteresis on two stiffness measures,
\begin{equation}
  \zeta = \max_p \frac{\Delta t}{t_{s,p}}, \qquad
  \chi  = \max \frac{\Delta t\,\rho_d}{\rho_g\,t_s},
\end{equation}
entering split-BE when either exceeds \code{hybrid\_enter\_\{zeta,chi\}} ($0.5$) and
returning to PC2 only when both fall below \code{hybrid\_exit\_\{zeta,chi\}} ($0.25$);
\code{hybrid\_force\_mode = pc2 | split\_be} pins a branch, and \code{coupling = pc2}
is that pin made permanent. The choice is reduced across ranks so every rank takes the
same branch, and a \code{DUST\_HYBRID\_SUMMARY} line reports the cycle counts and the
last $\zeta$, $\chi$. In the task graph (\file{dust/dust\_tasks.cpp}) this appears as
per-stage gates on every exchange and two migration chains; the deposited $Q,\mathbf P$
field gains a fifth feedback-rate channel; a second timestep task follows the split-BE
migration.

\paragraph{Consequence for the program.} The science targets ($\mathrm{St} = 0.1$--$10$
at $\Delta t \sim 10^{-2}\,\Om^{-1}$) sit in the resolved regime $\zeta \ll 1$, so
\code{coupling = pc2} under \code{rk2} is the natural production setting: second
order, about twice the speed of the IMEX dc1 path by the fork's measurement, no
\code{imex2+} CFL penalty, and the same integrator as all gravity validation. Note
that \code{hybrid} with the three-species NSH test ($t_{s,\min} = 0.01$) already
trips $\zeta = 0.55$ and runs split-BE throughout; the validation document records
the first-order signature that produces. The IMEX path stays as the reference for
stiff configurations.

\subsection{Build notes}

\code{imex2+} needed nothing: it is upstream \#612 and was already in the tree; the
dust fork only refactored its coefficients into \fn{SetImEx2PlusCoefficients}. Dust
sources are part of every build (\file{src/CMakeLists.txt}); the four dust problem
generators are built-in (\code{-D PROBLEM=built\_in\_pgens}, selected by
\code{<problem>/pgen\_name}). The rehearsal clone was built without network by
linking the tree's \file{kokkos} checkout and pointing
\code{FETCHCONTENT\_SOURCE\_DIR\_KOKKOS-FFT} at the cached \file{build/\_deps/kokkos-fft-src}.

%=========================================================================================
\section{Phase 4b (i): the shear-periodic deposit fold}
\label{sec:phase4b1}

\subsection{The problem}

A particle near a MeshBlock boundary deposits part of its weight into ghost cells of
its own block; the additive exchange (\S\ref{sec:rankpacked}) adds those ghost deposits
into the active cells of the neighbor that owns the location. Across a shear-periodic
$x_1$ face the owner's cells sit at a \emph{different} $y$: the tree wraps
\code{shear\_periodic} exactly like \code{periodic} (\file{mesh/meshblock\_tree.cpp}), so
the $x_1$ buffers of the face blocks carry the deposits to the unshifted $y$ --- the
``exact only for azimuthally uniform states'' caveat the module used to print. The copy
exchanges of hydro and of $u^\ast$ fix this by \emph{overwriting} the $x_1$ ghosts with
a remapped copy after the plain exchange; an additive exchange cannot overwrite, so the
unsheared contribution has to be suppressed and replaced.

\subsection{The sign}

The shearing-box identification used throughout the tree is
$u(x, y) = u(x + L_x, y - \qom t)$: the copy exchange fills the inner-$x_1$ ghosts with
the outer boundary's content shifted by $+y_{\rm sh}$ ($y_{\rm sh} = \qom L_x t$), and
the outer ghosts by $-y_{\rm sh}$ (\file{multigrid/multigrid\_shear.cpp:144--146},
\file{gravity/fft\_gravity.cpp:414--416}). A \emph{deposit} in an inner ghost cell at
$y$ is a contribution to the outer boundary at $y - y_{\rm sh}$; the deposit map is the
adjoint (transpose) of the ghost-fill map, so its signs are reversed:

\invariant{Adjoint signs.}{Inner-ghost deposits are shifted by $-y_{\rm sh}$ and
added to the \emph{outer} face blocks; outer-ghost deposits by $+y_{\rm sh}$ into the
\emph{inner} face blocks. Copying the multigrid fill's per-face signs verbatim gives
the wrong answer.}

\subsection{The fold}

\fn{MeshBoundaryValuesDep::FoldShearDeposit(a, yshear, rcon)} (\file{bvals/bvals\_dep.cpp})
reuses the Phase-1 global-plane machinery of \fn{Multigrid::FillShearGhostPack} with
three changes: the gather reads \emph{ghost} slabs and is additive, the signs are the
adjoint ones, and the scatter \emph{adds} into active cells.
\begin{enumerate}\itemsep2pt
  \item \textbf{Gather.} Two planes \code{shear\_plane\_(face, v*ng+d, gk, gj)} of the
        global $y$--$z$ extent (uniform grid: \code{mesh\_indcs.nx2}$\times$\code{nx3}),
        zeroed, then every face block adds its whole $x_1$ ghost slab --- all $k, j$
        \emph{including the $y$/$z$ corner ghosts} --- at the periodically wrapped
        global index (face 0 $\leftarrow$ inner ghosts \code{is-1-d}, face 1
        $\leftarrow$ outer ghosts \code{ie+1+d}). The corner rows overlap the active
        rows of the $y$/$z$ neighbors' slabs, so the gather is an atomic
        \emph{sum}: unlike the multigrid fill, several blocks legitimately write the same
        plane cell. Rows beyond a non-periodic $z$ face are dropped (the vertical-boundary
        policy of \S\ref{sec:next} item 7; the dust guards currently require periodic $z$).
  \item \textbf{MPI seam.} \code{Kokkos::fence} and one \code{MPI\_Allreduce(SUM)} of each
        plane, reached by every rank in lockstep (the fold is a task; all ranks run the
        same graph, and every send a rank waits for is posted before the collective).
  \item \textbf{Remap.} Per (face, variable, depth, $z$-row): the integer part of the
        shift is folded into a periodically wrapped scratch load, the fractional part
        \code{eps = fmod(ysh, dy)/dy} goes through \fn{DC/PLM/PPMX\_RemapFlx}
        (\file{shearing\_box/remap\_fluxes.hpp}) and the conservative update
        $q_j \leftarrow q_j - (F_{j+1} - F_j)$; $dy = L_y/\code{gny}$ is the global
        rank-consistent cell width (the \fn{ConsistentDx2} lesson). Face 0 uses
        $-y_{\rm sh}$, face 1 $+y_{\rm sh}$.
  \item \textbf{Scatter-add.} Face 0 into the outer face blocks' \code{ie-d}, face 1
        into the inner blocks' \code{is+d}, active $k, j$ only.
\end{enumerate}
The remap preserves the row sum, so the total deposited mass is conserved to
round-off; the summation order in the gather and the reduction can differ with the
rank count at round-off level only.

\subsection{Suppressing the unsheared contribution}

\fn{RecvAndSumDeposit} skips, on a block whose inner (outer) $x_1$ face is
shear-periodic, every receive buffer whose direction has $o_{x1} < 0$ ($> 0$): faces,
$x_1x_2$ and $x_3x_1$ edges, and corners --- all of which carry $x_1$-ghost content of
the neighbor across the face. The direction of buffer index $n$ comes from a 56-entry
table built in the constructor from the slot layout documented in
\file{mesh/nghbr\_index.hpp} (faces 0--7, $x_1x_2$ edges 16--23, $x_3x_1$ edges 32--39,
corners 48--55) and cross-checked against \fn{NeighborIndex} at construction.

\paragraph{A pitfall worth recording.} The first version built that table by calling
\fn{NeighborIndex} over all direction and sub-index arguments. For edges the function
uses only the first sub-index, and the second landed on \emph{other} slots --- the
$x_3x_1$ edges of one sign overwrote those of the other. The plain pass then still
delivered the unsheared $x_3x_1$-corner deposits, which the unit test caught as an
off-by-one in the multiplicity check (\S\ref{sec:phase4b1test}) and as a $10^{-4}$ drift
of the 3D NSH equilibrium.

\subsection{Hooks and parameters}

Two tasks, \fn{DustGasDrag::FoldDepQP} and \fn{FoldPMBR} (\file{dust/dust\_tasks.cpp}),
follow \fn{RecvDepQP} and \fn{RecvPMBR} with the same per-coupling stage gates; the
implicit solve now depends on the fold. \fn{CompleteAddExchange} (coupled solvers) calls
the fold as well. The offset is \fn{DustGasDrag::ShearOffset()} $= \qom L_x\,t$ at
\code{pmesh->time} on every stage, the convention of the $u^\ast$ remap.
\code{<dust>/shear\_remap} selects the order (\code{plm} default, as for $u^\ast$;
\code{ppmx} available), \code{<dust>/shear\_fold = false} is a diagnostic that skips the
fold (the face blocks' $x_1$ deposits are then dropped, not unsheared). The warning the
module printed under shear is replaced by an information line.

\subsection{The test generator}
\label{sec:phase4b1test}

\code{pgen\_name = dust\_deposit\_shear} (\file{pgen/tests/dust\_deposit\_shear.cpp},
\file{inputs/dust/dust\_deposit\_shear.athinput}) exercises the exchange without
particles: the $x_1$ ghost slabs of the face blocks are filled with known data, every
other cell is zero, and the exchange runs exactly as the module runs it. Ghost deposits
are \emph{partial} sums (a block's corner ghosts overlap its neighbors' active rows),
which the first version of the test got wrong by filling every ghost cell with the full
function; the corrected test carries three fields: a smooth $g(x,y,z)$ periodic in $y,z$
on the slab rows inside the block's own $y$/$z$ range (one source per location, so the
strips must equal $g$ at the shifted location up to the remap error), a constant on the
same rows (exact for any shift), and a constant on the \emph{whole} slab whose received
value must equal the analytic overlap multiplicity $(1 + n_y)(1 + n_z)$ of the source
row (exact at integer shifts). The generator also reports the largest value leaked into
any active cell outside the receiving strips (must be exactly zero) and the relative
mass change (round-off). \code{problem/fold = false} runs the plain pass alone.

\subsection{A structured, rank-independent initial condition}

The NSH generator's random placement hashes each particle's tag to a block index
\emph{within the local pack}, so the realization depends on the rank count (as does
\code{dust\_damping}'s, seeded by the pack's first GID) --- fine for physics runs,
useless for serial-vs-MPI comparisons; a divergence that this produced was first
mistaken for a bug in the fold (the no-back-reaction control agreed to $10^{-13}$
because test particles never feed positions back). \code{dust\_nsh} therefore gained
\code{<problem>/mass\_mod\_amp}, \code{mass\_mod\_phase}: on the lattice, particle
masses are multiplied by $1 + A\sin(2\pi (y - y_{\min})/L_y + \phi)$, a deterministic
azimuthally structured dust distribution (with \code{check\_equilibrium = false}).
\file{inputs/dust/dust\_nsh\_shear3d.athinput} is the tracked 3D shear-periodic NSH
input with all Phase-4b keys present for command-line overrides.

%=========================================================================================
\section{Phase 4b (ii): the ideal-gas energy update}
\label{sec:phase4b2}

\subsection{The problem}

The dust fork applied the drag back-reaction to the gas momentum only and therefore
required an isothermal EOS (a fatal guard). Gravito-turbulence needs cooling, hence an
ideal gas, and then every change of the gas momentum must be accompanied by a change of
the total energy: otherwise the internal energy $e = E - |\mathbf M|^2/2\rho$ silently
absorbs the negative of the kinetic-energy change.

\subsection{The rule}

\fn{GasKick(u0, m,k,j,i, dm, ideal)} (\file{dust/dust.hpp}) is the single entry point
for a drag kick $\Delta\mathbf M$ on a cell: for an ideal gas it adds
$\Delta E = \big(|\mathbf M + \Delta\mathbf M|^2 - |\mathbf M|^2\big)/2\rho$ before
adding $\Delta\mathbf M$, so the kick leaves $e$ unchanged. The four sites where the
module changes the gas momentum all use it: the IMEX back-reaction apply
(\fn{ApplyPMBR}), the IMEX history term (\fn{AddDragHistoryGas},
$\Delta\mathbf M = \tilde a_{22}\Delta t\,\mathbf R_g$), the PC2 predictor apply
(\fn{GasImplicitSolve}, stage 1) and the hybrid commit (\fn{ApplyPMBR}, stage 2, with
the tableau-weighted predictor removal folded into the same kick). The coupled solvers
only read the gas state and apply through \fn{ApplyPMBR}.

\paragraph{Frictional heating (\code{<dust>/drag\_heating}, default \code{false}).}
The dissipation of one drag kick $\Delta\mathbf v_j = c_j(\tilde{\mathbf u}_j - \mathbf v_j)$
on particle $j$ --- the kinetic energy lost by particle and gas together beyond the gas
kinetic-energy change that \fn{GasKick} books --- is
\begin{equation}
  Q_j = m_j\,c_j\,(1 - c_j/2)\,|\tilde{\mathbf u}_j - \mathbf v_j|^2 \;\ge 0 ,
\end{equation}
deposited with the particle's weights into a fourth channel of the back-reaction field
\code{dmom} (now four variables; the exchange objects follow) and added to $E$ where the
momentum channels are applied; the IMEX history term scales it with the same
$\tilde a_{22}$. On the PC2 path the midpoint impulse gives
$Q_j = m_j[\Delta\mathbf v_{\rm drag}\cdot(\mathbf u_g - \mathbf v_h)
- \tfrac12|\Delta\mathbf v_{\rm drag}|^2]$, the midpoint quadrature of
$m_j|\mathbf u - \mathbf v|^2/t_s$ over $\Delta t$ in kinetic-energy form. With an
isothermal EOS the flag is ignored with a warning.

\subsection{What the bookkeeping cannot do, and the lesson from the first test}

Inside a low-storage RK stage the conserved state is combined linearly,
$\mathbf u \leftarrow \gamma_0\mathbf u + \gamma_1\mathbf u_1$; $E$ combines linearly
but $|\mathbf M|^2/2\rho$ does not, so the combined internal energy exceeds the
combination of internal energies by $\gamma_0\gamma_1|\mathbf M_a - \mathbf M_b|^2/2\rho$
per step, where $\mathbf M_a - \mathbf M_b$ is the stage-1 momentum change --- for a
drag-dominated cell, the stage-1 kick. The default mode therefore does \emph{not} keep
$e$ exactly fixed over a step; it gains this positive gap, first order in $\Delta t$
(the kick scales with $\Delta t$ while $\Delta t \ll t_s$) and sub-first-order while
$\Delta t \gtrsim t_s$ of the stiffest species, where the kick saturates. The particle
velocities are combined the same way, so the heating mode carries the analogous gap on
the dust side: exact conservation of $E_{\rm gas} + K_{\rm dust}$ is not attainable in
this framework, and the measured residuals (validation document \S6) are this
truncation error. Both modes converge with $\Delta t$; both are inert for an
isothermal gas.

The first energy test looked like a 30\% bookkeeping error; it was the history file's
six significant digits (a total energy of $2.5$ cannot show changes of $10^{-6}$).
\file{inputs/dust/dust\_damping\_ideal.athinput} therefore uses $p_{\rm gas} = 10^{-5}$,
comparable to the dissipated energy, and \code{<time>/dtmax} to pin the step. The NSH
generator gained the ideal-gas initialization it lacked (uniform \code{<problem>/pgas});
\file{inputs/dust/dust\_nsh\_shear3d\_ideal.athinput} is the ideal-gas twin of the 3D
input.

%=========================================================================================
\section{Phase 4b (iii): the particle timestep under shear}
\label{sec:phase4b3}

\subsection{The problem}

The inherited transport limit (\S\ref{sec:dt}) is $\Delta t \le c_{\rm p}\,\min_p
\Delta x_2/|v_y - \qom x|$ with $c_{\rm p} = $ \code{dt\_cfl} $= 0.5$: one cell of the
\emph{total} azimuthal motion per step. The background shear is a smooth translation
the drift integrates exactly ($y \leftarrow y + \Delta t\,(v_y - \qom x)$ is linear in
$x$), not a signal the mesh has to resolve; the gas, orbitally advected, is limited by
its peculiar velocity only. In the $L_x = 1$ test boxes the two limits happen to
coincide with the sound-speed limit, which is why the penalty never showed. In an
$L_x = 8$ box with the NSH input's $\Delta x_2 = 1/32$ the particle limit is
$0.5\cdot(1/32)/(1.5\cdot 4) = 2.6\times10^{-3}$ against a gas limit of
$1.37\times10^{-2}$: a factor $5.3$, growing with $L_x$.

\subsection{What the limit is actually for}

Two mechanisms depend on how far a particle moves in one stage, and neither is a
cell:

\begin{enumerate}\itemsep3pt
\item \textbf{The deposit halo.} The TSC stencil reaches one cell beyond the particle,
      and \code{MeshBoundaryValuesDep} exchanges one ghost layer. Both are satisfied
      whenever the particle lies \emph{inside its owning block} at deposit time ---
      which is what the per-stage migration establishes: after
      \fn{ParticlesBoundaryValues::SetNewPrtclGID} every particle carries the
      \code{gid} of the block that contains it. The halo-violation abort in
      \fn{DepositDrag} (\file{dust/dust\_pm.cpp}) is precisely the detector of a failed
      migration.
\item \textbf{The migration.} \fn{SetNewPrtclGID} forms the integer offset of the
      particle from its block, $i_x = \lfloor (x - x_{\min} + L_x^{\rm mb})/L_x^{\rm mb}
      \rfloor - 1$ and likewise in $y$, $z$, and indexes the neighbor table with
      \fn{NeighborIndex}$(i_x,i_y,i_z,\ldots)$. The table covers faces, edges and
      corners: offsets in $\{-1,0,1\}$. A particle that traverses a whole MeshBlock in
      one stage produces $|i_y| = 2$, an out-of-range neighbor slot, and a wrong or
      invalid destination. The one exception is a crossing of the shear-periodic
      $x_1$ mesh boundary, routed by the $(y,z)$ shear maps of Phase 4a for any
      azimuthal offset.
\end{enumerate}

So the hard requirement is \emph{less than one MeshBlock per stage}, in every
direction, on the total transport velocity; the cell-per-step limit on the peculiar
velocity stays for accuracy, as for the gas.

\subsection{The rule}

\fn{DustGasDrag::ComputeNewTimeStep} (\file{dust/dust\_newdt.cpp}) now reduces two
minima over the particles in one kernel and takes the smaller:
\begin{align}
\Delta t_{\rm cell}  &= c_{\rm p}\,\min_p \min_i \frac{\Delta x_i}{|v_i^{\rm cell}|},
   & v_y^{\rm cell} &= \begin{cases} v_y & \text{\code{dt\_transport = relative}}\\
                                      v_y - \qom x & \text{\code{full}} \end{cases}\\
\Delta t_{\rm block} &= \frac{s}{\beta_{\max}}\,\min_p \min_i
   \frac{L_i^{\rm mb}}{|v_i^{\rm full}|},
   & \beta_{\max} &= \max_{\rm stages} |\beta_s|,
\end{align}
with $v^{\rm full}_y = v_y - \qom x$ in the 3D shearing box (and $v^{\rm full} = v$
otherwise), $L_i^{\rm mb}$ the block extent, and $s = $ \code{<dust>/dt\_block\_safety}
(default $0.5$). The factor $\beta_{\max}$ is what makes $s$ the literal fraction of a
block crossed per stage. In a low-storage stage the position is combined as
$y \leftarrow \gamma_0 y + \gamma_1 y_1 + \beta_s\Delta t\,\dot y$ with
$\gamma_0 + \gamma_1 = 1$, so the displacement from the position the stage started
with is $\gamma_1(y_1 - y) + \beta_s\Delta t\,\dot y$; for the tableaus in use the
register term never adds to the drift term in magnitude, and the per-stage
displacement is bounded by $\beta_{\max}\Delta t\,|\dot y|$:

\begin{center}
\small
\begin{tabular}{@{}lccc@{}}
\toprule
integrator & stage displacements $/\,(\Delta t\,|\dot y|)$ & $\beta_{\max}$ & bound used \\
\midrule
\code{imex2+}, $\gamma = 1 + 1/\sqrt2$ & $1.707,\ 0.707$ & $1.707$ & $1.707$ \\
\code{imex2+}, $\gamma = 1/2$ (switch) & $0.5,\ 0.5$ & $1.0$ & $1.0$ \\
\code{rk2} & $1.0,\ 0$ & $1.0$ & $1.0$ \\
\code{rk3} & $1.0,\ 0.5,\ 0.5$ & $1.0$ & $1.0$ \\
\bottomrule
\end{tabular}
\end{center}

The kernel reads \code{pdrive->beta[]} over \code{nexp\_stages}, so the $\gamma$-switch
is covered without a table. The premise $\gamma_0 + \gamma_1 = 1$ with non-negative
weights holds for every tableau the module's integrator guard admits (\code{imex2+} for
\code{coupling = imex}, \code{rk2} for \code{pc2}/\code{hybrid}) and for \code{rk3};
\code{rk4} has $\gamma_0 < 0$ and would need its own bound, but cannot be selected. Note the pre-existing asymmetry the
cell limit carries under \code{imex2+}: $c_{\rm p} = 0.5$ is applied to $\Delta t$, while
stage 1 advances $\gamma\Delta t = 1.71\Delta t$, i.e.\ $0.85$ cells per stage; it is
unchanged here and harmless (nothing in the module needs less than a block).

\subsection{Parameters and diagnostics}

\code{<dust>/dt\_transport = relative|full} (default \code{relative}; \code{full} is
the legacy limit, bit-identical to Phase 4b (ii)) and \code{<dust>/dt\_block\_safety}
(default $0.5$; a value $\ge 1$ is accepted with a warning for the diagnostic run of
validation \S7 that shows the halo abort the guard prevents; the constructor prints
which limit is in force in a 3D shearing box). The guard is always on, in all three
directions and in every geometry; outside a shearing box it is inert in practice
(the cell limit is stricter by $N_i^{\rm mb}$). The destructor prints
\code{DUST\_DT\_SUMMARY block\_guard\_cycles=$n$}, the number of cycles in which the
block guard was below the cell limit \emph{on rank 0} --- a rank-local count that says
which of the two particle limits was the active one, not whether the particles limited
the run at all (the gas usually does once the shear is out). The two 3D
shear-periodic inputs carry both keys so they can be varied from the command line
(\S\ref{sec:params}).

\subsection{What the larger step exposes}

Once the shear is out of the limit, the step in a wide box is the gas step, and for
species with $t_s$ below it the IMEX coupling runs in its stiff regime. The validation
scan (\S7 there) shows this is a plain, converging truncation error --- second order,
$1\%$ in the quiet phase and $26\%$ during a streaming burst at $\Delta t/t_{s,\min}
= 1.4$ --- that the legacy limit had been masking in wide boxes only. The operating
rule that follows is recorded in the validation document; it does not change the
timestep code, which now limits what must be limited and nothing else.

\subsection{What does not change}

The push, the deposits, the fold and the migration are untouched: a larger step simply
means that particles now hop MeshBlocks routinely in $y$ (up to $s\,N_2^{\rm mb}$ cells
per stage) and cross the shear-periodic face with multi-cell azimuthal offsets, both
paths that existed and were exercised before at one cell per step. The split-BE branch
of the hybrid coupling keeps its later timestep task (\fn{NewTimeStep2}), which calls
the same routine.

%=========================================================================================
\section{What the next phases must add}
\label{sec:next}

The survey established these gaps; they are listed in the order they bite.

\begin{enumerate}\itemsep4pt
\item \textbf{Shear-periodic additive deposit remap (4b).} \textbf{Done},
      \S\ref{sec:phase4b1}. The Phase-4c gravity deposit will be a third additive field
      and inherits the fold through \code{MeshBoundaryValuesDep}. The alternative
      design (image particles deposited directly in the neighbor's frame) stays in
      reserve; the field remap is exact for constants and second/third order otherwise,
      and it costs one small collective per fold.
\item \textbf{Ideal-gas back-reaction (4b).} \textbf{Done}, \S\ref{sec:phase4b2}.
\item \textbf{Particle timestep under shear (4b).} \textbf{Done},
      \S\ref{sec:phase4b3}: limit on $v_y$, with a MeshBlock guard on the total
      velocity.
\item \textbf{Dust self-gravity (4c).} A plain mass deposit $\sum_j m_j W_{jk}/V$ into
      a new \code{rho\_dust} (with the additive exchange of item 1), added at the three
      reads of \S\ref{sec:map}; the $\mathbf g$ field, its exchange, and the gather in
      \fn{ExplicitPush}; skip the dormant-stage solve. The FFT solver is extended too so
      it remains the oracle (decided).
\item \textbf{Particle restart (4d).} \file{outputs/restart.cpp} ignores the particle
      arrays; long cluster runs are restart chains.
\item \textbf{Diagnostics (4d).} \code{dust\_d} is an NGP binning while the physics
      deposit is TSC; \code{pvtk} omits stopping time and mass; the SC14 history's
      gravitational stress becomes a \emph{total}-gravity stress once dust contributes
      to $\Phi$ and needs its own dust columns; a Roche-density / Hill-radius clump
      finder does not exist.
\item \textbf{Vertical boundary policy (4c/4d).} Particles reaching the open (diode)
      $z$ faces and deposits landing in their ghost zones: the Athena-C reference
      silently discards both; decide explicitly and count the mass.
\item \textbf{Refinement (Phase 5).} Dust is fatal on \code{multilevel} in three
      places, and upstream particles have no remesh handling at all (\code{PGID} goes
      stale after the first regrid; nothing in \file{mesh\_refinement.cpp} touches
      particles). Needs the deposit exchange across levels and a particle
      redistribution hook next to \fn{ShearingBox::ReinitAfterMeshUpdate}.
\item \textbf{Upstream defect.} \code{Particles::dtnew} is never assigned on the
      cosmic-ray path yet is read by \fn{Mesh::NewTimeStep}; the dust module sets it,
      plain particle runs do not.
\end{enumerate}

%=========================================================================================
\section{Input-file interface}
\label{sec:params}

Required: \code{<time>/integrator = imex2+} for \code{coupling = imex}, \code{rk2}
for \code{pc2} or \code{hybrid}; \code{<particles>/particle\_type = dust},
\code{pusher = imex\_dust}, \code{ppc} (particles per cell, an integer multiple of
\code{nspecies} for lattice placement); \code{<mesh>/nghost >= 2}; periodic
boundaries, or shear-periodic $x_1$ in 3D. Any EOS (the isothermal requirement for back-reaction is gone, \S\ref{sec:phase4b2}). Shearing-box runs read \code{<shearing\_box>/qshear},
\code{omega0}, \code{stratified}.

\begin{center}
\small
\begin{tabular}{@{}>{\raggedright\arraybackslash}p{4.3cm}>{\raggedright\arraybackslash}p{2.5cm}>{\raggedright\arraybackslash}p{8.6cm}@{}}
\toprule
\code{<dust>} parameter & default & meaning \\
\midrule
\code{coupling} & imex & \code{imex} (needs \code{imex2+}) / \code{hybrid} / \code{pc2}
  (both need \code{rk2}); see \S\ref{sec:merge2} \\
\code{shear\_remap} & plm & \code{dc} / \code{plm} / \code{ppmx}: $y$-remap order of the
  shear-periodic deposit fold (\S\ref{sec:phase4b1}) \\
\code{shear\_fold} & true & \code{false}: diagnostic, skip the fold \\
\code{drag\_heating} & false & ideal gas: deposit the frictional dissipation of each
  drag kick as heat (\S\ref{sec:phase4b2}); default keeps $e$ fixed and discards it \\
\code{hybrid\_force\_mode} & auto & \code{auto} / \code{pc2} / \code{split\_be}
  (\code{hybrid} only) \\
\code{hybrid\_enter\_zeta}, \code{hybrid\_enter\_chi} & 0.5, 0.5 & split-BE entry
  thresholds on $\zeta$, $\chi$ \\
\code{hybrid\_exit\_zeta}, \code{hybrid\_exit\_chi} & 0.25, 0.25 & PC2 re-entry thresholds \\
\code{nspecies} & 1 & number of species; \code{taus\_1..N} (positive) required \\
\code{dust\_to\_gas} & 0.01 & total dust-to-gas ratio for the default mass initializer \\
\code{back\_reaction} & true & deposit the drag reaction onto the gas \\
\code{deposit} & tsc & \code{ngp} / \code{cic} / \code{tsc} assignment and gather kernel \\
\code{dt\_cfl} & 0.5 & particle transport CFL factor (cell-crossing limit) \\
\code{dt\_transport} & relative & \code{relative}: the cell limit uses the peculiar
  $v_y$; \code{full}: the legacy $v_y - \qom x$ (\S\ref{sec:phase4b3}) \\
\code{dt\_block\_safety} & 0.5 & fraction of a MeshBlock a particle may cross per
  stage on the total velocity ($\ge 1$: guard off, warns; diagnostic) \\
\code{gamma\_switch} & false & Krapp et al.\ eq.~18: $\gamma \to 1/2$ when
  $\Delta t > \max t_s$ \\
\code{stopping\_time\_mode} & species\_fixed & or \code{particle\_static}, \code{dynamic} \\
\code{drag\_solver} & local & \code{local} / \code{dc1} / \code{dc2} / \code{pcg} /
  \code{adaptive} / \code{applya} (coupled Schur-complement solvers; \code{local} is
  the BLKP19 closed form and the validated baseline) \\
\code{drag\_rtol}, \code{drag\_atol}, \code{drag\_iter\_max} & $10^{-11}$, $10^{-14}$, 200 &
  PCG tolerances \\
\code{drag\_adaptive\_*} & --- & error target of the adaptive solver \\
\code{drag\_diagnostic\_interval} & 0 & \code{DUST\_SOLVER\_DIAG} print cadence \\
\bottomrule
\end{tabular}
\end{center}

Problem-generator parameters (\code{rho0}, \code{etavk}, \code{eps\_s}, placement,
seeds, eigenmode amplitudes; \code{dust\_nsh}'s \code{mass\_mod\_amp},
\code{mass\_mod\_phase}, \code{check\_equilibrium}; \code{dust\_deposit\_shear}'s
\code{time0}, \code{remap}, \code{fold}) are specific to the five built-in generators
\code{dust\_damping}, \code{dust\_nsh}, \code{dusty\_wave}, \code{streaming\_linear},
\code{dust\_deposit\_shear};
the files under \file{inputs/dust/} and \file{../athenak\_dust\_tests/inputs/} are
the reference configurations. Outputs: \code{file\_type = pvtk} with
\code{variable = prtcl\_all} (positions, gid, tag, species, velocities);
\code{variable = dust\_d} on any mesh output (NGP mass density, needs a \code{<dust>}
block); \code{<time>/dtmax} caps the step for fixed-step convergence tests.

\end{document}
